\chapter{局部紧交换群}

在本章中，除非另有说明，字母 G 表示一个赋有 Haar 测度（通常记作 dx）的局部紧交换群；对 \(p \in [1, +\infty]\)，空间 \(L^{p}(G, dx)\) 将简记为 \(L^{p}(G)\)，其范数记作 \(f \mapsto \|f\|_{p}\)。我们把 \(L^{1}(G)\) 等同于 \(\mathcal{M}^{1}(G)\) 的一个子空间，等同映射为 \(f \mapsto f \cdot dx\)。回忆 Haar 测度的支集等于 G（INT，VII，§1，第 1 节，注 3）；特别地（INT，III，§2，第 2 节，命题 9），从空间 \(\mathcal{K}(G)\) 到 \(L^{p}(G)\) 的典范映射对 \(p \in [1, +\infty]\) 是单射。对 \(p \neq +\infty\)，我们把 \(L^{p}(G)\) 等同于空间 \(\widetilde{\mathcal{F}}(G; \mathbf{C})\) 的一个子空间，后者是 G 上几乎处处有定义且取有限值的复值函数的类所成的空间（INT，IV，§3，第 5、6 节）。特别地，记号 \(L^{1}(G) \cap L^{2}(G)\) 表示 \(L^{1}(G)\) 与 \(L^{2}(G)\) 在该空间中的交。我们用 \(f \mapsto \widetilde{f}\) 表示对合代数 \(L^{1}(G)\) 上的对合（I，第 99 页，例 4）；\(\widetilde{f}(x) = \overline{f(x^{-1})}\) 对 G 中一切 \(x\) 成立。

回忆（INT，V，§5，第 3 节，定理 1）：若 μ 是局部紧拓扑空间 X 上的一个复测度，f 是 X 上一个局部 μ-可积函数，则以 f 为密度关于 μ 的测度 ν 记作 f·μ 或简记作 fμ。从 X 到 C 的函数 g 关于测度 ν 本质可积当且仅当 gf 关于 \(\mu\) 本质可积；此时有

\[
(f \cdot \mu) (g) = \int_ {\mathrm{X}} g d (f \cdot \mu) = \int_ {\mathrm{X}} g f d \mu = \mu (g f).
\]

若 G 是紧拓扑群，G 上的规范化 Haar 测度是 G 上唯一的 Haar 测度 \(\mu\)，它满足 \(\mu(\mathrm{G}) = 1\)。

若 \(X\) 是离散拓扑空间，\(X\) 上的计数测度是离散测度 \(\mu\)（在 \(X\) 上），它满足 \(\mu(\{x\}) = 1\) 对所有 \(x \in X\) 成立。若 \(G\) 是离散拓扑群，则 \(G\) 上的计数测度是 \(G\) 上的一个 Haar 测度。

我们还将用到下面两个引理。

引理 1.　设 X 为局部紧拓扑空间，\(\mu\) 为 X 上的一个正测度。设 \(x\) 为 \(\mu\) 的支集中的一个元素。设 U 为 \(x\) 的一个开邻域。存在一个支集含于 U 的函数 \(f \in \mathcal{K}_{+}(\mathrm{X})\)，使得 \(\int fd\mu = 1\)。

由测度支集的定义（INT，III，§2，第 2 节，定义 1），存在一个支集含于 U 的函数 \(g \in \mathcal{K}(\mathrm{X})\)，使得 \(\mu(g) \neq 0\)。我们有 \(\mu(|g|) > 0\)，因为 \(\mu\) 是正测度，而函数 \(f = \mu(|g|)^{-1}|g|\) 即具有所需的性质。

引理 2.　设 G 与 H 为分别以 \(e_{\mathrm{G}}\) 与 \(e_{\mathrm{H}}\) 为单位元的拓扑群。假设拓扑群 G 是 Hausdorff 的。设 f 为从 G 到 H 的一个拓扑群态射。假设对 G 中 \(e_{\mathrm{G}}\) 的每个邻域 U，存在 H 中 \(e_{\mathrm{H}}\) 的一个邻域 W，使得 \(f^{-1}(\mathrm{W}) \subset \mathrm{U}\)。则态射 f 是单的且是严格的（TG，III，第 16 页，定义 1）。

这些假设蕴含 \(\operatorname{Ker}(f)\) 包含于 G 中 \(e_{G}\) 的每个邻域，从而由于 G 是分离的，有 \(\operatorname{Ker}(f) = \{e_{\mathrm{G}}\}\)。于是 f 通过过渡到子空间所诱导的从 G 到 \(f(\mathrm{G})\) 的同态是双射；用 \(g: f(\mathrm{G}) \to \mathrm{G}\) 表示逆同态。这些假设进而蕴含 g 在 \(e_{H}\) 处连续，从而连续（TG，III，第 15 页，命题 23）。

